%**********************************************************
%                1. CLASSE E OPÇÃO DO DOCUMENTO
%
% Ative UMA das opções abaixo conforme o tipo de documento.
% Mantenha as demais comentadas.
% NÃO modifique o arquivo classPPGMSB.cls.
%**********************************************************

% --- Documentos em PORTUGUÊS ---
\documentclass[diss]{classPPGMSB}   % Dissertação de mestrado  ← PADRÃO
%\documentclass[tese]{classPPGMSB}  % Tese de doutorado
%\documentclass[qlf]{classPPGMSB}   % Exame de qualificação

% --- Documentos em INGLÊS ---
%\documentclass[diss=en]{classPPGMSB}  % Master's dissertation
%\documentclass[tese=en]{classPPGMSB}  % Doctoral thesis
% Nota: [diss,ing] é sintaxe equivalente a [diss=en] (legada)


%***************************************************
%                     PACOTES INICIAIS
%***************************************************
\usepackage{lipsum}

% O pacote hyperref já é carregado pela classe classPPGMSB.
% Use apenas \hypersetup para configurar as cores dos links.
\hypersetup{
	colorlinks = true,
	linkcolor = ppgmsbC1,
	citecolor = ppgmsbC1,
	urlcolor = ppgmsbC1
}
%=====================================================
     %     ELEMENTOS PARA CRIAÇÃO DE SIGLAS
%           CRIE AQUI SUA LISTA DE SIGLAS
%=====================================================
% A classe já carrega o pacote glossaries-extra e o estilo das listas.
% NÃO use \usepackage{glossaries-extra}, \makeglossaries nem \printglossary:
% as listas são geradas só pelo LaTeX (pdflatex/xelatex/lualatex ou latexmk),
% sem makeglossaries e sem nenhum passo extra de compilação.
%
% Defina símbolos e abreviações nos arquivos abaixo, sempre aqui no preâmbulo.
% Só aparecem nas listas as entradas usadas no texto com \gls{...}.
%================ ACESSE ESSES ARQUIVOS =============
%*************** INCLUA AS SIGLAS E ABREVIAÇÕES =====
% Incluir os arquivos de siglas e abreviações
% Definir símbolos
\newglossaryentry{pi}{
	name={\ensuremath{\pi}}, 
	symbol={\ensuremath{\pi}}, % Adiciona o símbolo explicitamente
	description={A razão entre a circunferência de um círculo e seu diâmetro}, 
	type=symbols
}

\newglossaryentry{alpha}{
	name={\ensuremath{\alpha}}, 
	symbol={\ensuremath{\alpha}}, % Adiciona o símbolo explicitamente
	description={Constante de rotação em mecânica clássica}, 
	type=symbols
}

\newglossaryentry{epsilon}{
	name={\ensuremath{\varepsilon}}, 
	symbol={\ensuremath{\varepsilon}}, % Adiciona o símbolo explicitamente
	description={Pequena variação ou tolerância em um sistema}, 
	type=symbols
}

\newglossaryentry{lambda}{
	name={\ensuremath{\lambda}}, 
	symbol={\ensuremath{\lambda}}, % Adiciona o símbolo explicitamente
	description={Comprimento de onda na teoria das ondas}, 
	type=symbols
}

\newglossaryentry{gamma}{
	name={\ensuremath{\gamma}}, 
	symbol={\ensuremath{\gamma}}, % Adiciona o símbolo explicitamente
	description={Fator de Lorentz na relatividade especial}, 
	type=symbols
}

% Definir abreviações
\newabbreviation{ai}{AI}{Inteligência Artificial}
\newabbreviation{iot}{IoT}{Internet das Coisas}
\newabbreviation{ml}{ML}{Aprendizado de Máquina}
\newabbreviation{nasa}{NASA}{National Aeronautics and Space Administration}
\newabbreviation{www}{WWW}{World Wide Web}
\newabbreviation{api}{API}{Application Programming Interface}


%====INSTRUÇÕES DE USO=========
% Usando uma abreviação no texto pela primeira vez
% A primeira vez que \gls{ai} é usado, aparece "Inteligência Artificial (AI)"

% Usando uma abreviação no texto novamente
% Uso subsequente de \gls{ai} mostra apenas "AI"

% Usando um símbolo no texto
% Usando \gls{pi} no texto mostra o símbolo \pi com sua descrição na lista de símbolos"

%********************************************************
%======== COMANDO INTERNO DA CLASSE===============
\opcoesDocumento
% ================================================

\begin{document}
%=====================================================
%     INÍCIO DOS DADOS GERAIS DO DOCUMENTO
%=====================================================


%=====================================================
% INFORMAÇÕES AUTOR E DA DISSERTAÇÃO
%=====================================================

%========== INFORMAÇÕES DA INSTITUIÇÃO =================
\Universidade{Universidade do Estado da Bahia}
\Departamento{Departamento de Ciências Exatas e da Terra --- Campus II}
\Programa{Programa de Pós-Graduação em Modelagem e Simulação de Biossistemas}
\UrlUNEB{https://www.uneb.br}
\UrlMSB{https://ppgmsb.uneb.br}

%========== LOGOS DA INSTITUIÇÃO ======================
% Coloque os arquivos de imagem na pasta indicada.
% Baixe os logos em: https://ppgmsb.uneb.br
% Se o arquivo não existir, um quadrado com X é exibido no lugar.

\LogoUniversidade{figuras/logo_uneb.jpg}
\LogoPrograma{figuras/ppgmsbLogo.png}
\LogoCabecalho{figuras/logoso.png}

% Dados da Gestão da Universidade
\ReitorNome{Adriana Marmori Lima}
\ProreitorNome{Elizeu Clementino de Souza}
\DiretorNome{Issac Douglas Moreira}
\CoordenadorNome{Marcos Batista Figueredo}

%========== DADOS DO AUTOR DO TRABALHO ====================

%========== NOME COMPLETO SEM ABREVIAÇÕES =================
% --- Autor(a) ---
\NomeAutor{Juliana Ribeiro Santos}
\EmailAutor{discente.completo@ppgmsb.test}
\AutorNacionalidade{Brasileiro(a)}
\AutorMunicipio{Alagoinhas}
\AutorEstado{BA}
\LinkLattes{http://lattes.cnpq.br/8000000000000008}
\LinkOrcid{0000-0002-1825-0097}
\DataNascimento{12/04/1996}

%============= TITULO DA DISSERTAÇÃO =================
%---- Título não deve execeder 150 caracteres --------
%========TITULO EM PORTUGUES
\TituloBR{Modelagem Computacional e Simulação Estocástica de Biossistemas Aplicadas à Vigilância Epidemiológica}
%=========TÍTULO EM INGLES==========
\TituloEN{Computational Modeling and Stochastic Simulation of Biosystems Applied to Epidemiological Surveillance}
%=====================SUBTITULO===================
%-----------Caso não tenha comente a linha -----
%----------- limitado a 80 caracteres ----------
%\SubTituloBR{limitado a 80 caracteres}
%\SubTituloEN{subtitle within the 80-character limit.}

%=====================SUBTITULO===================
% Palavras-chave (PT e EN) — entre 3 e 7 termos, separados por ";"
\palavrasChavePT{modelagem computacional; simulação estocástica; biossistemas; aprendizado de máquina; vigilância epidemiológica}
\palavrasChaveEN{computational modeling; stochastic simulation; biosystems; machine learning; epidemiological surveillance}
\palavrasChaveBusca{modelagem computacional; simulação estocástica; vigilância epidemiológica; UNEB; PPGMSB}




%===================== ORIENTADOR DO TRABALHO ===================
\OrientadorNome{Beatriz Carvalho Mendes}
\OrientadorFiliacao{Universidade do Estado da Bahia}
\OrientadorSigla{UNEB}
\OrientadorLattes{http://lattes.cnpq.br/9000000000000003}
\OrientadorOrcid{0000-0002-9000-0003}
%===================== ODS DO TRABALHO ===================
\odsSelecionadas{3,7,11,13,15}

%===================== DADOS DE ADERÊNCIA DO TRABALHO 
\AreaDeConcentracaoBR{Área de Teste}
\AreaDeConcentracaoEN{Test Area}
\LinhaPesquisaBR{Linha de Teste}
\LinhaPesquisaEN{Test Research Line}
\GrupoDePesquisaBR{Grupo de Modelagem e Simulação de Biossistemas}
\GrupoDePesquisaEN{Biosystems Modeling and Simulation Group}
%===================== DADOS ADICIONAIS ===============

\DataDefesa{30/11/2026}
\HoraDefesa{14:00h}
\LocalDefesa{Alagoinhas/BA}

%========== LICENÇA (NOVO) =================
\Licenca[https://creativecommons.org/licenses/by-nc-sa/4.0/]{CC BY-NC-SA 4.0}

%========== FINANCIAMENTO (NOVO) =================
% Agências de fomento, bolsas, projetos
\Financiamento{%
	CAPES --- Coordenação de Aperfeiçoamento de Pessoal de Nível Superior (Bolsa de Mestrado)
}

%========== REPOSITÓRIO DE CÓDIGO (NOVO) =================
% Se houver código fonte disponível
\Repositorio{https://github.com/usuario/repositorio}

%========== VOLUMES (NOVO) =================
% Número de volumes do trabalho (padrão: 1)
\Volumes{1}

%=============== INFORMADO PELA SECRETARIA ============
\UrlTrabalho{https://doi.org/xxxxxxxx}     % se houver página/handle/doi
\Msbdoi{10.1234/ppgmsb.xxxxx}              % se aplicável

%========================================================
%      INDICAÇÃO DOS COORIENTADORES
%========================================================

% ================== COORIENTADORES (A até C) ==================
% A — obrigatório só se existir coorientador
\CoorientadorANome{Rafael Teixeira Nunes}
\CoorientadorAFiliacao{Universidade do Estado da Bahia}
\CoorientadorASigla{UNEB}
\CoorientadorAFormacao{Ciência da Computação}
\CoorientadorAAtuacao{Aprendizado de Máquina e Modelagem Epidemiológica}
\CoorientadorAOrcid{0000-0002-9000-0004}
\CoorientadorALattes{http://lattes.cnpq.br/9000000000000004}

% B — se houver
%\CoorientadorBNome{Nome do Coorientador B}
%\CoorientadorBFiliacao{Instituição do Coorientador B}
%\CoorientadorBSigla{SIGLA-B}
%\CoorientadorBFormacao{Área de formação B}
%\CoorientadorBAtuacao{Área de atuação B}
%\CoorientadorBOrcid{0000-0000-0000-0000}
%\CoorientadorBLattes{http://lattes.cnpq.br/BBBBBBBBBBBBBBBB}

% C — se houver
%\CoorientadorCNome{Nome do Coorientador C}
%\CoorientadorCFiliacao{Instituição do Coorientador C}
%\CoorientadorCSigla{SIGLA-C}
%\CoorientadorCFormacao{Área de formação C}
%\CoorientadorCAtuacao{Área de atuação C}
%\CoorientadorCOrcid{0000-0000-0000-0000}
%\CoorientadorCLattes{http://lattes.cnpq.br/CCCCCCCCCCCCCCCC}

%*******************************************************
%      MEMBROS DA BANCA DA DEFESA E/OU QUALIFICAÇÃO
%*******************************************************
%
%Indique o numero de membros da banca considerando: 
% Maximo 5 Para mestrado/Doutorado
% Mínimo 3 Para mestrado (Orientador, Membro Interno e Externo)

%Indique 1 para interno
%indique 2 para externo
%Lembre que o orientador sempre será membro da banca
%========================================
\numeroMembrosBanca{3}
%=========================================
% ================== MEMBROS DA BANCA EXAMINADORA ==================

% Tipo: 1 = Membro Interno | outro valor = Membro Externo

% --- Membro A ---
\MembroANome{Luciana Prado Freitas}
\MembroAFiliacao{Universidade do Estado da Bahia}
\MembroASigla{UNEB}
\MembroAFormacao{Matemática Aplicada}
\MembroAAtuacao{Modelagem Matemática em Saúde}
\MembroAOrcid{0000-0002-9000-0005}
\MembroALattes{http://lattes.cnpq.br/9000000000000005}
\MembroATipo{1}

% --- Membro B ---
\MembroBNome{Eduardo Moura Bastos}
\MembroBFiliacao{Universidade Federal da Bahia}
\MembroBSigla{UFBA}
\MembroBFormacao{Estatística}
\MembroBAtuacao{Inferência Bayesiana e Processos Estocásticos}
\MembroBOrcid{0000-0002-9000-0006}
\MembroBLattes{http://lattes.cnpq.br/9000000000000006}
\MembroBTipo{2}

% --- Membro C ---
%\MembroCNome{Profa. Ana Paula Menezes}
%\MembroCFiliacao{Universidade Estadual de Feira de Santana}
%\MembroCSigla{UEFS}
%\MembroCFormacao{Matemática}
%\MembroCAtuacao{Métodos Numéricos e Inteligência Artificial}
%\MembroCOrcid{0000-0001-6358-2274}
%\MembroCLattes{http://lattes.cnpq.br/3333333333333333}
%\MembroCTipo{1}

% --- Membro D ---
%\MembroDNome{Prof. Ricardo Tavares de Oliveira}
%\MembroDFiliacao{Instituto Federal da Bahia}
%\MembroDSigla{IFBA}
%\MembroDFormacao{Ciência da Computação}
%\MembroDAtuacao{Aprendizado de Máquina e Processamento de Sinais}
%\MembroDOrcid{0000-0002-1944-5566}
%\MembroDLattes{http://lattes.cnpq.br/4444444444444444}
%\MembroDTipo{2}

% --- Membro E ---
%\MembroENome{Profa. Mariana Castro Nogueira}
%\MembroEFiliacao{Universidade Federal do Recôncavo da Bahia}
%\MembroESigla{UFRB}
%\MembroEFormacao{Engenharia Ambiental}
%\MembroEAtuacao{Sistemas Complexos e Modelagem Hidrológica}
%\MembroEOrcid{0000-0003-6701-8472}
%\MembroELattes{http://lattes.cnpq.br/5555555555555555}
%\MembroETipo{2}

%=====================================================


%=====================================================
%     INÍCIO DO PRÉ-TEXTO
%
% Durante o processo de escrita recomendamos que os elementos pré-textuais fiquem comentados, o que acelera bastante o processo de compilação
%=====================================================
%===================CAPA============================
%              GERAÇÃO AUTOMÁTICA 
%=====================================================
\makeCapa
%===================================================

%===================CONTRA-CAPA============================
%              GERAÇÃO AUTOMÁTICA 
\makeContraCapa
%===================================================

%===================FOLHA APROVAÇÃO================
%              GERAÇÃO AUTOMÁTICA 
\folhaAprovacao
%===================================================
%\pretextual
\frontmatter
%===================AGRADECIMENTOS================
\begin{agradecimentos}
	\lipsum[1-2]
\end{agradecimentos}

%===================================================

%===================DEDICATÓRIA================
\begin{dedicar}
	\lipsum[3][1-4]
\end{dedicar}

%===================================================

%===================RESUMO================
% COMENTAR CASO DA VERSÃO EM INGLES
%===================RESUMO================
% INSERIR AQUI SEU RESUMO E 
\begin{resumo}
	A modelagem computacional e a simulação estocástica de biossistemas desempenham papel fundamental na compreensão de fenômenos epidemiológicos complexos e no suporte à tomada de decisão em saúde pública. Neste trabalho fictício de teste para o sistema de gestão de bancas do PPGMSB, propomos uma arquitetura híbrida que integra equações diferenciais estocásticas e algoritmos de aprendizado de máquina para a predição espaço-temporal de surtos epidemiológicos em cenários com dados ruidosos. O método compreende três etapas principais: pré-processamento e filtragem adaptativa de séries temporais, calibração paramétrica via inferência bayesiana e simulação de cenários de intervenção. Os experimentos numéricos sintéticos demonstram estabilidade numérica, redução do erro quadrático médio de predição e elevada aderência aos critérios de reprodutibilidade científica do programa.
\end{resumo}
% 
%===================================================

%=================== ABSTRACT ==============

\begin{newAbstract}
	Computational modeling and stochastic simulation of biosystems play a fundamental role in understanding complex epidemiological phenomena and in supporting decision-making in public health. In this fictitious test work for the PPGMSB defense management system, we propose a hybrid architecture that combines stochastic differential equations and machine learning algorithms for the spatio-temporal prediction of epidemiological outbreaks in noisy data scenarios. The method comprises three main stages: preprocessing and adaptive filtering of time series, parametric calibration through Bayesian inference, and simulation of intervention scenarios. Synthetic numerical experiments show numerical stability, a reduction of the mean squared prediction error, and high adherence to the scientific reproducibility criteria of the program.
\end{newAbstract}
%===================================================

%===================listas =========================
%              GERAÇÃO AUTOMÁTICA 
%===================================================

% Lista de Símbolos (não é impressa se nenhum símbolo for usado no texto)
\listaDeSimbolos
% Lista de Abreviações (não é impressa se nenhuma abreviação for usada no texto)
\listaDeAbreviacoes
%lista de Figuras
\listaDeIlustracoes
%Lista de Tabelas
\listaDeTabelas
%Lista de Algoritmos
\listaDeAlgoritmos

% Desativar a numeração de páginas para a tabela de conteúdos
\pagestyle{empty} % Remove a numeração da página atual

% Sumário automático
\setlength{\beforechapskip}{-50pt}
\tableofcontents*

% Reiniciar a numeração de páginas e voltar ao estilo normal
\clearpage % Finalizar a página atual
\pagestyle{plain} % Estilo de página com numeração
\setcounter{page}{1} % Reiniciar a numeração de páginas

%**********************************************************
%                FIM ELEMENTOS PRÉ-TEXTUAIS
%**********************************************************

\mainmatter
\pagestyle{ppgmsbPages}
%**********************************************************
%                INÍCIO DOS ELEMENTOS TEXTUAIS
%*******************************************************
%===================================================

%====================== CAPITULOS =====================
% Para redigir seu documento abra os arquivos e modifique o texto
% Você pode modificar os elmentos textuais conforme orientação. Mas observe sempre os critérios mínimos.
\chapter{Introdução}
\label{cap:introducao}

\lipsum[1]

O uso de \gls{ml} e de \gls{ai} na vigilância epidemiológica tem crescido \cite{Suppa2022,Lahmiri2019}. \lipsum[2][1-3]

\section{Contextualização}
\lipsum[2-3]

\section{Objetivos}
\lipsum[3-4]

\subsection{Objetivo geral}
\lipsum[6]

\subsection{Objetivos específicos}
\lipsum[7]

\section{Organização do trabalho}
\lipsum[4-5]

\chapter{Revisão de literatura}
\label{cap:revisao}

\lipsum[1]

Trabalhos recentes \cite{Almeida2019,Nassif2022,Rieder2020} discutem a taxa de transmissão $\gls{lambda}$ e o parâmetro $\gls{gamma}$. \lipsum[3][1-3]

\section{Modelos compartimentais}
\lipsum[2-3]

\subsection{Modelos determinísticos}
\lipsum[5]

\subsection{Modelos estocásticos}
\lipsum[6]

\section{Aprendizado de máquina em saúde}
\lipsum[3-4]

\chapter{Materiais e métodos}
\label{cap:modelo}

\lipsum[1]

Seja $\gls{epsilon}$ o ruído de observação e $\gls{alpha}$ a taxa de incubação \cite{Andreis2022}. Os dados foram obtidos por \gls{api} pública.

\section{Dados}
\lipsum[2-3]

\section{Modelo proposto}
\lipsum[3-4]

\subsection{Calibração bayesiana}
\lipsum[6]

\subsection{Simulação de cenários}
\lipsum[7]

\section{Métricas de avaliação}
\lipsum[4-5]

\chapter{Resultados e discussão}
\label{cap:resultados}

\lipsum[1]

Os resultados são comparados com \cite{Saba2022}.

\section{Cenário base}
\lipsum[2-3]

\section{Cenários de intervenção}
\lipsum[3-4]

\section{Discussão}
\lipsum[4-5]

\chapter{Conclusão}
\label{cap:conclusao}

\lipsum[1]

\section{Trabalhos futuros}
\lipsum[2-3]





%====================== FIM DOS CAPITULOS =====================

%============ELEMENTOS PÓS-TEXTUAIS==========

%%----------------------------------------------------------------------------
%% Configurar as referencias bibliograficas
%%----------------------------------------------------------------------------
\clearpage % Garante que a bibliografia comece em uma nova página
%\addcontentsline{toc}{chapter}{Referências}

\bibliography{posTextual/referencias}

%Voce pode alterar o arquivo de referencias a vontade desde que coloque na pasta corretamente
%Para citações use:% Exemplo de citação no texto

%\cite{<<autor>>} -> para citação direta.

%\citep{<<autor>>} -> para citação indireta.

%Para figuras use o comando \fonte
%% Exemplo de uso do comando \fonte
%\begin{figure}[h]
%	\centering
%    \caption{Exemplo de figura}
%	\includegraphics[width=0.5\textwidth]{figura_exemplo}
%	\fonte{Próprio autor} ou \fonte{\citep{<<autor>>}}
%\end{figure}

% Para as ODS, use \odsSelecionadas no bloco de metadados:
% \odsSelecionadas{3,7,13,15}
% A página de metadados chama \ODSFiguras automaticamente.




%==========================================
% POS TEXTUAL APENDICES
%==========================================
%\begin{appendices}
%	\chapter{Questionário}
Quisque facilisis auctor sapien. Pellentesque gravida hendrerit lectus. Mauris rutrum
sodales sapien. Fusce hendrerit sem vel lorem. Integer pellentesque massa vel augue. Integer
elit tortor, feugiat quis, sagittis et, ornare non, lacus. Vestibulum posuere pellentesque eros.
Quisque venenatis ipsum dictum nulla. Aliquam quis quam non metus eleifend interdum.
Nam eget sapien ac mauris malesuada adipiscing. Etiam eleifend neque sed quam. Nulla
facilisi. Proin a ligula. Sed id dui eu nibh egestas tincidunt. Suspendisse arcu.
%\end{appendices}

%==========================================
%\renewcommand{\restoreapp}{}
%\renewcommand\appendixname{ANEXOS}
%\renewcommand\appendixpagename{ANEXOS}
%\begin{appendices}
%	\chapter{Vantagens}
A grande vantagem do ambiente appendices é que, uma vez que o ambiente seja terminado, você pode continuar com seções ou capítulos como antes — a numeração não é afetada pelos anexos intervenientes.
O pacote fornece outra maneira alternativa de definir anexos, como divisões inferiores no documento. O ambiente subappendices permite que você coloque anexos separados para uma seção específica, codificados como subsections, ou para um capítulo específico, codificados como sections. Então, alguém poderia escrever:
%\end{appendices}
\paginaMetadados

\end{document}

%=====================================================

